\documentclass[11pt,a4paper]{article}
\usepackage[a4paper,margin=22mm]{geometry}
\usepackage{fontspec,babel}
\babelprovide[main,import]{russian}
\babelfont{rm}{Times New Roman}
\babelfont{sf}{Arial}
\setmonofont{Consolas}
\usepackage{amsmath,amssymb,booktabs,graphicx,xcolor,caption,hyperref}
\hypersetup{colorlinks,allcolors=blue!50!black,pdftitle={Эффективная динамика идеальной 12-электродной ловушки}}
\captionsetup{font=small,labelfont=bf,hypcap=false}
\setlength{\parindent}{0pt}
\setlength{\parskip}{4pt}
\setlength{\emergencystretch}{2em}
\newcommand{\dd}{\mathrm{d}}
\newcommand{\R}{\mathbf R}
\newcommand{\V}{\dot{\mathbf R}}
\newcommand{\Lref}{L}
\newcommand{\e}[1]{\times10^{#1}}
\begin{document}
{\sffamily\small ИССЛЕДОВАТЕЛЬСКИЙ ОТЧЁТ \hfill Идеальная ловушка / версия 1}\par
\vspace{0.55cm}
{\sffamily\LARGE\bfseries Восстановление эффективной динамики\\[2pt]
трёхфазной 12-электродной ловушки\par}
\vspace{0.25cm}
{\small 8 ноября 2026 г. \quad $q=10^{-12}$ Кл, $m=10^{-15}$ кг, $U_0=250$ В, $f=100$ кГц}

\section*{Ответ на физический вопрос}
Периоды при нескольких амплитудах позволяют восстановить ведущую радиальную
энергию в заданном степенном семействе. В синтетическом контроле восстанавливается
$p=6$. Единственный проверенный радиус старта в FEM определяет нормировку при
фиксированном $p$, но не саму степень. Подгонка этой нормировки сохраняет точный
период и почти постоянный радиус разворота, однако оставляет поперечную ошибку.

Поправка к ускорению $\gamma(r)J\V$ с $\gamma(r)\propto r^4$ выведена из двух
RF-квадратур и проверена без подгонки коэффициента к траекториям.
На отложенной RF-фазе она снижает векторную RMS-ошибку с $0{,}556\%$ до
$0{,}00104\%$ от 1 мм. В аналитическом полном поле эффект сохраняется,
меняет знак при обращении вращения и наблюдается при противоположных
скоростях в практически одной медленной точке.

\section*{Данные, область проверки и протокол}
\begin{tabular}{@{}p{.28\linewidth}p{.67\linewidth}@{}}
\toprule
Набор & Использование в версии 1 \\
\midrule
P4 h35, фаза $0$ & 10,25 $T_s$, единственный запуск для оценки $B_6^\ast$.
Короткий префикс того же запуска не считается независимым. \\
P4 h35, фаза $\pi/2$ & 3,25 $T_s$, отложенное полное начальное состояние.
Номинальный медленный старт у обоих запусков $(1\text{ мм},0)$, $V_0=0$. \\
P4 h50 и шаги & Те же старты для оценки численной чувствительности.
В v2 проверены $T_{\rm RF}/50$ и $T_{\rm RF}/100$, для пилота фазы 0
также $T_{\rm RF}/200$ и независимый Radau. \\
Пилоты P2 и N48 & P2 и старые проверки края/перезапуска дают историю численных
ограничений. Диагностические A\textsubscript{fit} на N48 не являются прямыми FEM-данными.
Из N48 берётся поле модели B и карта микродвижения. \\
50 $T_s$ & Сохранённый расчёт h35, $T_{\rm RF}/100$ от $t=0$.
Используется только опубликованная диагностика, около $13^\circ$.
Глобальная сходимость этого интервала ещё не измерена. \\
\bottomrule
\end{tabular}

В исходной проверке v2 межсеточные RMS h50/h35 равны $0{,}09337\%$ (фаза 0)
и $0{,}03239\%$ (фаза $\pi/2$). Численные допуски v2 по RMS, максимуму
и периоду равны $0{,}1/0{,}3/0{,}1\%$. Физические пороги сравнения
сохранены, $1/3/1\%$. SHA исходных массивов, журналов переходов и родительских
массивов сверены. Данные FEM использованы из готовых контрольных точек.

\newpage
\section*{1. Обратная задача и синтетический контроль}
Для радиального старта из покоя, при $\Psi=U_L(r/L)^p$, $L=1$ мм, и амплитуде
$a$ из интеграла энергии получается
\[
T(a)=4L\sqrt{\frac{m}{2U_L}}\,
\frac{\mathrm B(1/p,1/2)}{p}\left(\frac{a}{L}\right)^{1-p/2}.
\]
Наклон $\log T$ по $\log(a/L)$ определяет $p$. Свободный член даёт $U_L$,
а коэффициент в СИ равен $B=U_L/L^p$. При одном радиусе матрица регрессии
имеет ранг 1. Измерения на нескольких радиусах дают ранг 2 в степенном
семействе. Для произвольной формы потенциала конечное число периодов не
определяет функцию. В непрерывной постановке с известным $T(E)$ обращение
Абеля требует дополнительных условий на ветви потенциала~\cite{tautochrone}.

Период оценён по временам максимумов $r$ через один максимум. Это полный
период знакопеременного колебания. Локальное уточнение использует кубический
сплайн $\R(t)$. Амплитуда есть средний радиус этих точек. Угол линии
определён как $\tfrac12\operatorname{unwrap}(2\operatorname{atan2}(y,x))$.
Из наклона угла по физическому времени получается скорость поворота.
Отдельно записан спектр проекции на первую ось разворота за один измеренный
полный цикл. Гармоники служат диагностикой формы, угол определяется по
фазе радиальных разворотов.

\begin{center}
\begin{tabular}{@{}lrrr@{}}
\toprule
Выборка C & $a$, мм & $T$, мс & $1/T$, Гц \\
\midrule
Обучение & 0,50 & 92,365324 & 10,826574 \\
Обучение & 1,00 & 23,091331 & 43,306295 \\
Обучение & 1,25 & 14,778452 & 67,666086 \\
Проверка & 0,75 & 41,051255 & 24,359791 \\
Проверка & 0,90 & 28,507816 & 35,078099 \\
Проверка & 1,10 & 19,083745 & 52,400617 \\
\bottomrule
\end{tabular}
\end{center}
При сингулярных числах $1{,}757$ и $0{,}666$ получено $p-6=1{,}15\e{-13}$.
Относительная ошибка $U_L$ составляет $-7{,}08\e{-10}$, как и ошибка
численного значения $B$ в выбранных единицах СИ с этой точностью.
Максимальная ошибка предсказанного периода на отложенных амплитудах
составляет $3{,}55\e{-10}$ в относительных единицах, а координатная RMS
меньше $4{,}1\e{-7}\%$. Эти числа характеризуют данный синтетический расчёт.

\begin{center}
\includegraphics[width=\linewidth]{potential.png}
\captionof{figure}{Восстановленные энергии на общей шкале. Правая панель показывает
малую разность с C. Продолжение кривой D за 1 мм использует предположение $p=6$.}
\end{center}

\newpage
\section*{2. Скалярная модель и чётность по скорости}
В идеальной геометрии поворот на $\pi/6$ эквивалентен сдвигу RF-фазы
на $-2\pi/3$ по модулю периода. При записи $\sin(\Omega t-n\theta+\delta_n)$
разрешены знаковые порядки $n\equiv4\pmod{12}$, начиная с $4,-8,16$.
Стационарный скаляр инвариантен под поворотом на $\pi/6$. Для
отражения относительно оси электрода RF-уравнение инвариантно при
одновременном обращении времени со сдвигом RF-фазы. В медленных
координатах, где отражение и обращение времени действуют обычным образом,
это даёт скаляру угловые члены $\cos(12k\theta)$ при выборе оси $x$
вдоль зеркальной оси. Первая допустимая аналитическая степень угловой
поправки равна 12, но в ведущем псевдопотенциале первый вклад
от гармоник 4 и 16 имеет вид $r^{18}\cos(12\theta)$.
Квадратичный вклад от $-8$ даёт $r^{14}$. Перекрёстный вклад 4 и $-8$
сокращается в сумме квадратов градиентов двух квадратур.

Оценка по полю даёт $B_{14}=5{,}73\e{17}\ \text{Дж}\,\text{м}^{-14}$, или
$2{,}59\e{-8}$ от $B_6r^6$ при $r=1$ мм. Аналогичная оценка первой
анизотропии даёт $2{,}53\e{-13}$ от $B_6r^6$ на том же радиусе.
Две RF-фазы при одном радиусе не обеспечивают амплитудного рычага
для оценки $B_{14}^\ast$. В D подогнан только $B_6^\ast$.
Радиальная модель сохраняет $L_z=x\dot y-y\dot x$. При строго
радиальном старте она сохраняет и исходную прямую. Малый $L_z$
после обратной карты учтён одинаково во всех моделях.

\subsection*{Вывод скоростного члена}
Пусть $\ddot{\mathbf r}=a(\mathbf r,\chi)$, где
$a=f_c\cos\chi+f_s\sin\chi$ и $\chi=\Omega t+\phi_0$.
В калибровке $\mathbf r=\R+\xi$ с нулевым быстрым средним $\xi$ разложение
до первого члена, линейного по $\V$, даёт
\[
\xi=-\frac{a}{\Omega^2}
      -\frac{2}{\Omega^3}D(\partial_\chi a)\,\V+\cdots.
\]
Второй член компенсирует смешанную производную $2\Omega\partial_\chi
(\V\cdot\nabla)\xi$. Усреднение $Da\,\xi$ даёт градиент ведущего
псевдопотенциала и скоростное ускорение
\[
-\Omega^{-3}[Df_c,Df_s],\V,\qquad [X,Y]=XY-YX.
\]
Симметрия якобианов электростатических сил делает коммутатор
кососимметричным. Для заданного направления октуполя получаем
$\gamma(r)J\V$, где $J(x,y)=(-y,x)$ и
\[
\gamma(r)=\frac{288q^2C_4^2}{m^2\Omega^3}r^4,
\qquad
\gamma(L)=1{,}26760747\ \text{с}^{-1}.
\]
Обращение последовательности фаз меняет знак. Эта асимптотика
выводится для выбранных медленных координат. Синтетические силы при
быстром периодическом возбуждении возникают и в иных схемах
усреднения~\cite{goldman,higashikawa}.

Член не меняет энергию $m|\V|^2/2+B_6r^6$, поскольку
$\V\cdot J\V=0$. В чистом октуполе сохраняется
$L_z-\gamma(L)r^6/(6L^4)$. Наклон в этой зависимости даёт независимую
оценку коэффициента по сглаженным $\R,\V$. В E оставлено аналитическое
$B_6$ из C, чтобы измерить вклад скоростной поправки отдельно от подгонки D.

\newpage
\section*{3. Прямой P4 FEM и отложенная RF-фаза}
Оценка A использует центрированное среднее координат за $T_{\rm RF}$
на узлах через $T_{\rm RF}$. Контроли используют сырой строб на
фиксированной RF-фазе, обратную карту с поправкой на дрейф, а также
двукратное среднее по $2T_{\rm RF}$ для момента импульса.
Начальные $(\R,\V)$ для B--E получены одной обратной картой из полного
состояния при $t=0$. Используется общий физический таймер без
фазового и пространственного совмещения. Поэтому результат B для
фазы $\pi/2$ немного отличается от прежнего сравнения с номинальным
стартом. Ошибка угла измерена в моменты разворотов A, по модулю $180^\circ$.

\begin{center}
\small
\begin{tabular}{@{}llrrrrr@{}}
\toprule
Запуск & Модель & RMS, \% & Max, \% & $|\Delta T/T|$, \% & RMS$_\perp$, \% & RMS угла, $^\circ$ \\
\midrule
$0$, 10,25 $T_s$ & B & 1,7490 & 4,5005 & $9{,}30\e{-5}$ & 1,7487 & 1,5446 \\
& C & 1,7490 & 4,5005 & $9{,}47\e{-5}$ & 1,7487 & 1,5446 \\
& D & 1,7490 & 4,5005 & $1{,}11\e{-4}$ & 1,7487 & 1,5446 \\
& E & 0,00130 & 0,00252 & $4{,}82\e{-5}$ & 0,00013 & 0,00011 \\
\midrule
$\pi/2$, 3,25 $T_s$ & B & 0,55607 & 1,35246 & $9{,}99\e{-5}$ & 0,55606 & 0,50297 \\
& C & 0,55607 & 1,35246 & $1{,}02\e{-4}$ & 0,55606 & 0,50297 \\
& D & 0,55608 & 1,35246 & $1{,}04\e{-4}$ & 0,55606 & 0,50297 \\
& E & 0,00103 & 0,00232 & $4{,}84\e{-5}$ & 0,00096 & 0,00086 \\
\bottomrule
\end{tabular}
\end{center}

Оба периода A составляют около $23{,}091309$ мс. Радиусы разворота
в среднем $0{,}99999945$ и $0{,}99999941$ мм. Подгонка D даёт
$B_6^\ast=22{,}12401492\ \text{Дж}\,\text{м}^{-6}$ против
$22{,}12392402\ \text{Дж}\,\text{м}^{-6}$ для C. По трём основным фильтрам
диапазон D составляет $[22{,}123966;22{,}126752]$, включая
сырой строб, сохраняющий фазовый сдвиг микродвижения.
Перенос той же оценки на h50 меняет $B_6^\ast$ примерно на $7{,}4\e{-5}$
в относительных единицах, больше сдвига D относительно C ($4{,}1\e{-6}$).
В фазе 0 наблюдается $+11{,}16553^\circ$/с, в фазе $\pi/2$
$+11{,}99607^\circ$/с. Для E это $+11{,}16473$ и $+11{,}97693^\circ$/с.
Для B остаётся около $-0{,}0023$ и $+0{,}810^\circ$/с
из-за медленного начального момента импульса.

\begin{center}
\includegraphics[width=\linewidth]{trajectories.png}
\captionof{figure}{Поперечная координата и форма траектории. Кривые A и E
почти совпадают на этом масштабе. B, C и D практически совпадают между собой.}
\end{center}

\newpage
\section*{4. Проверки механизма и пределы вывода}
Для P4 изменение скорости поворота при переходе h35\,$\to$\,h50 равно
$0{,}00030^\circ$/с и $0{,}2753^\circ$/с соответственно. Переход
между шагами на h35 даёт менее $10^{-6}$ градуса в секунду в этих
двух сравнениях. Сырой строб меняет наклон на
$0{,}00140^\circ$/с и $0{,}00054^\circ$/с. Точность совпадения E
с h35 нельзя приписывать континуальному FEM-решению с той же погрешностью.
Межсеточная разница меньше улучшения E, а её вклад в абсолютную точность
E больше показанной ошибки относительно h35.

В независимом локальном контроле полного аналитического RF-поля
при $r=1$ мм использованы скорости $v_x=\pm0{,}12$ м/с. Интегрирование
в обе стороны от старта позволяет оценить медленный тренд в центре.
Совместная регрессия полинома и RF-гармоник на окнах $\pm32$ и
$\pm64$ RF-периодов при двух фазах дала разность $a_y=0{,}30422$--$0{,}30423$
м/с$^2$ при расхождении центральных положений до $4{,}5\e{-14}$ м.
Отношение к разности скоростей даёт
$\gamma(L)=1{,}267589$--$1{,}267621$ с$^{-1}$.
У локального скаляра сила в одной точке одинакова для противоположных
скоростей. Кроме того, для чистого вращающегося октуполя непрерывная
симметрия поворот плюс сдвиг фазы требует радиального стационарного
скаляра. Вместе с независимыми траекториями это исключает естественное
скалярное уравнение в выбранных медленных координатах на проверенной
точности. Утверждение относится к уравнению с обычной кинетической
энергией и силой, зависящей только от положения.

\begin{center}
\includegraphics[width=\linewidth]{analytic_control.png}
\captionof{figure}{Полная аналитическая RF-динамика при новых амплитудах и углах.
При обращении вращения отклик зеркально отражается. Скорость поворота
получена по разворотам, без FFT-оценки низкой частоты.}
\end{center}

\begin{center}
\small
\begin{tabular}{@{}lrrrr@{}}
\toprule
Старт, мм & RF, $^\circ$/с & E, $^\circ$/с & RMS C, \% & RMS E, \% \\
\midrule
0,75 & $+3{,}53253$ & $+3{,}53331$ & 0,23415 & 0,000121 \\
1,00 & $+11{,}16700$ & $+11{,}16701$ & 0,55515 & 0,000465 \\
1,25 & $+27{,}25259$ & $+27{,}26324$ & 1,08391 & 0,001480 \\
1,25, обратное вращение & $-27{,}25259$ & $-27{,}26324$ & 1,08391 & 0,001480 \\
\bottomrule
\end{tabular}
\end{center}
Расхождение при смене шага $T_{\rm RF}/10\to T_{\rm RF}/20$
не превышает $1{,}50\e{-7}\%$ RMS. Оценки $\gamma(L)$ по моменту
импульса на положительной последовательности составляют
$1{,}26729$--$1{,}26758$ с$^{-1}$ для чистого октуполя и
$1{,}26750$--$1{,}26883$ с$^{-1}$ для h35. Обратный запуск
при 1,25 мм зеркален по стартовому углу и служит контролем симметрии.

\newpage
\section*{5. Коэффициенты и следующий эксперимент}
\begin{center}
\begin{tabular}{@{}lll@{}}
\toprule
Источник & Параметр & Значение \\
\midrule
Поле P4 h35 & $C_4$ & $3{,}30419716\e{10}$ В\,м$^{-4}$ \\
Аналитическая C & $B_6$ & $22{,}1239240196$ Дж\,м$^{-6}$ \\
Обращение C & $B$ при $p\simeq 6$ & $22{,}1239240039$ в СИ \\
Подгонка D & $B_6^\ast$ & $22{,}1240149186$ Дж\,м$^{-6}$ \\
Поле, модель E & $\gamma(1\text{ мм})$ & $1{,}26760747$ с$^{-1}$ \\
Две скорости, чистое поле & $\gamma(1\text{ мм})$ & $1{,}267589$--$1{,}267621$ с$^{-1}$ \\
\bottomrule
\end{tabular}
\end{center}

\textbf{Проверено.} Восстанавливаемость параметров степенного семейства
на синтетике. Неудача радиальной скалярной модели в поперечном
канале P4. Количественное улучшение E на отложенной фазе и в чистом
RF-поле на новых радиусах. Аналитический контроль с двумя скоростями
проверяет нечётный по скорости вклад независимо от подгонки периода и угла.
Вдоль первой оси разворота для отложенной фазы A, B и E дают
$|X_3/X_1|\simeq0{,}068429$ и $|X_5/X_1|\simeq0{,}009473$.
Эти гармоники продольной формы почти нечувствительны к поперечному эффекту.

\textbf{Численные ограничения.} Вариация радиуса и азимута для прямого
FEM, вклад более высоких порядков и предел точности на 50 $T_s$
требуют отдельных проверок. Оценки чувствительности к шагу,
сетке и фильтру приведены раздельно. Они не являются статистическими
доверительными интервалами.

В прежней серии на гладком поле N48 одновременно умножали $U_0$ и $f$
на $s=1,2,4$. RMS составила соответственно $0{,}5552\%$, $0{,}2776\%$
и $0{,}1388\%$.
Это совместимо с $\gamma\propto s^{-1}$ при постоянном ведущем
потенциале. Серия не входила в подгонку и не измеряла угол напрямую.

\subsection*{Следующий эксперимент}
Предварительно зафиксировать E и прогнозы для прямого P4 FEM на
$0{,}75$ и $1{,}25$ мм при ненулевом азимуте и обоих направлениях
вращения. На 3,25 $T_s$ взять h50/h35, $T_{\rm RF}/50$ и
$T_{\rm RF}/100$ с одинаковыми полными стартами на каждой паре сеток.
Сравнить B и E по тем же порогам и проверить знаки и
межсеточную сходимость скоростей поворота. Затем, если межсеточный
разброс мешает проверке поправки, добавить h25. В чистом поле
дополнительно варьировать $f$ при фиксированных $U_0$ и старте,
проверяя закон $\gamma(L)\propto\Omega^{-3}$. Для 50 $T_s$ отдельно
сравнить наклоны по поворотам на независимых шагах и сетках
из $t=0$ перед количественным использованием длинной траектории.

\begin{thebibliography}{9}\small\raggedright
\bibitem{tautochrone} P. Terra, R. de Melo e Souza, C. Farina. Is the tautochrone curve unique?
\href{https://arxiv.org/abs/1610.01006}{arXiv:1610.01006}.
\bibitem{goldman} N. Goldman, J. Dalibard. Periodically driven quantum systems.
Phys. Rev. X 4, 031027 (2014).
\href{https://doi.org/10.1103/PhysRevX.4.031027}{doi:10.1103/PhysRevX.4.031027}.
\bibitem{higashikawa} S. Higashikawa, H. Fujita, M. Sato. Floquet engineering
of classical systems. \href{https://arxiv.org/abs/1810.01103}{arXiv:1810.01103}.
\end{thebibliography}
\end{document}
